% Shared Pandoc PDF code-block styling.
%
% Goals:
% - wrap long lines in both plain fenced blocks (`verbatim`) and syntax-highlighted blocks
% - keep a clear block container in the PDF for readability
% - avoid visible wrap markers that are not present in the source

\usepackage{fvextra}
\usepackage{framed}
\definecolor{codeblockbg}{RGB}{245,247,250}
\definecolor{codeblockborder}{RGB}{220,224,228}

\fvset{
  breaklines,
  breakanywhere,
  breaksymbolleft={},
  breaksymbolright={},
  breakindent=1em,
  fontsize=\small
}

% Pandoc emits untyped fenced code blocks as `verbatim`, so we recustomize it
% directly to keep long prompts/commands readable in PDFs.
\RecustomVerbatimEnvironment{verbatim}{Verbatim}{
  breaklines,
  breakanywhere,
  breaksymbolleft={},
  breaksymbolright={},
  breakindent=1em,
  fontsize=\small,
  frame=single,
  framesep=0.6em,
  rulecolor=\color{codeblockborder}
}

% Pandoc defines `Shaded` later in its template, so defer the redefinition
% until the document begins. This keeps syntax-highlighted blocks in a light
% container similar to the docs site while remaining print-friendly.
\makeatletter
\AtBeginDocument{
  \@ifundefined{Shaded}{
    \newenvironment{Shaded}{
      \def\FrameCommand{
        \fboxsep=0.6em
        \colorbox{codeblockbg}
      }
      \MakeFramed{\advance\hsize-\width\FrameRestore}
    }{
      \endMakeFramed
    }
  }{
    \renewenvironment{Shaded}{
      \def\FrameCommand{
        \fboxsep=0.6em
        \colorbox{codeblockbg}
      }
      \MakeFramed{\advance\hsize-\width\FrameRestore}
    }{
      \endMakeFramed
    }
  }
}
\makeatother
